\section{Finite-width transfer of the Gaussian covariance construction}
\label{sec:nongaussian-transfer}
Write $\phi=\tanh$. All comparison directions below are determined by
weights and auxiliary randomness, independently of the queried input.
The first layer has the Gaussian comparison law of
Lemma~\ref{lem:nongaussian-first-layer}; only the original prior is
non-Gaussian.

\begin{lemma}[Two-projection Lindeberg estimate]
\label{lem:two-projection-lindeberg}
Let $X=(X_1,\ldots,X_N)$ have independent coordinates of law $P$, let
$e$ be a unit vector, and let $b\in\mathbb R^N$ satisfy
$\|b\|_2\le1$ and $\|b\|_\infty\le N^{-1/2}$.
Put $r=e\cdot b$, $q=\|b\|_2^2$, and, for $W\ge1$ and $0<\nu\le1$,
\[
 a=\sqrt\nu\,b+(\sqrt W-\sqrt\nu)r e,
 \qquad \xi=e\cdot X,\qquad Z=a\cdot X.
\]
Let $(\xi_G,Z_G)$ be the centered Gaussian pair with the same
covariance. If $\theta=\|e\|_\infty$, then
\begin{align}
 |\E\phi(Z)^2-\E\phi(Z_G)^2|
 &\le C_P(1+W)^{3/2}(N^{-1/2}+\theta),\notag\\
 |\E[\xi\phi(Z)]-\E[\xi_G\phi(Z_G)]|
 &\le C_P(1+W)^{3/2}(N^{-1/2}+\theta).
 \label{eq:two-projection-lindeberg}
\end{align}
The covariance identities are exact:
\[
 \E\xi^2=1,\qquad \E Z^2=\nu q+(W-\nu)r^2,\qquad
 \E\xi Z=\sqrt W\,r.
\]
\end{lemma}
\begin{proof}
The covariance formulas use only centering and unit variance.
Moreover $\|a\|_2\le\sqrt W$ and
$\|a\|_\infty\le N^{-1/2}+\sqrt W\,\theta$.
Replace coordinates of $X$ one at a time by standard Gaussians.
The first two Taylor terms cancel because the first two moments
match. For $\phi(a\cdot X)^2$, the third derivative in coordinate
$j$ is bounded by $C|a_j|^3$. The total error is therefore at most
$C_P\sum_j|a_j|^3$.

For $(e\cdot X)\phi(a\cdot X)$, the third derivative is
\[
 (e\cdot X)a_j^3\phi'''(a\cdot X)
       +3e_j a_j^2\phi''(a\cdot X).
\]
During one replacement, separate
$e\cdot X=\xi_{-j}+e_jX_j$. The remaining mixed Gaussian/non-Gaussian
coordinates have covariance identity, so $\E|\xi_{-j}|\le1$.
Their independence from the replaced coordinate, and the fixed third
and fourth moments of the scalar laws, bound the Taylor remainder by
\[
 C_P\bigl(|a_j|^3+|e_j||a_j|^3+|e_j|a_j^2\bigr).
\]
Summing, using $\|a\|_2\le\sqrt W$ and the displayed infinity-norm
bound, proves \eqref{eq:two-projection-lindeberg}. The derivatives of
$t\mapsto\tanh t$ through order three are globally bounded.
\end{proof}

\begin{lemma}[A delocalized finite layer follows the Gaussian state map]
\label{lem:nongaussian-state-layer}
Condition on a previous hidden vector $h(x)$ and a unit direction $e$,
with $|h_i(x)|\le1$. Generate independent standardized rows
$X_i\sim P^N$ and use physical rows
\[
 \sqrt{\nu/N}\,T_{\sqrt{W/\nu}}X_i.
\]
Set $\xi_i=e\cdot X_i$, $e'_i=\xi_i/(\sum_j\xi_j^2)^{1/2}$,
$h'_i(x)=\phi(Z_i)$, and
\[
 q'=N^{-1}\sum_i h'_i(x)^2,\qquad
 r'=N^{-1/2}\sum_i e'_i h'_i(x),\qquad
 \sigma'=\sqrt{q'-(r')^2}.
\]
The zero denominator has probability zero. The associated Gaussian
state map is
\begin{align}
 \mathcal T_W(r,\sigma)
 &=\left(\E[\xi_GY],
      \|Y-\xi_G\E[\xi_GY]\|_{L^2}\right),\notag\\
 Y&=\phi(\sqrt W\,r\xi_G+\sqrt\nu\,\sigma\zeta_G),
 \label{eq:nongaussian-gaussian-state}
\end{align}
where $r=e\cdot h(x)/\sqrt N$, $\sigma^2=N^{-1}\|h(x)\|^2-r^2$,
and $\xi_G,\zeta_G$ are independent standard Gaussians.
For $0<u\le1/4$, except with probability $C e^{-c_PNu^2}$,
\begin{equation}
 \|(r',\sigma')-\mathcal T_W(r,\sigma)\|_2
 \le C_P\left[u+(1+W)^{3/2}
                 (N^{-1/2}+\|e\|_\infty)\right]^{1/2}.
 \label{eq:nongaussian-one-step-state}
\end{equation}
Also, with probability at least $1-\eta$,
\begin{equation}
 \|e'\|_\infty\le C_P\sqrt{\log(2N/\eta)/N},
 \label{eq:nongaussian-delocalization}
\end{equation}
provided $N\ge C_P\log(2/\eta)$.
\end{lemma}
\begin{proof}
Conditionally, the pairs $(\xi_i,Z_i)$ are IID. The variables
$\phi(Z_i)^2$ are bounded, $\xi_i\phi(Z_i)$ have sub-Gaussian tails,
and $\xi_i^2$ are subexponential with mean one. Standard scalar
Chernoff bounds show that their three empirical means differ from
expectation by at most $u$, with failure $C e^{-c_PNu^2}$.
In particular the normalizing denominator squared, divided by $N$,
lies in $[1-u,1+u]$.
Lemma~\ref{lem:two-projection-lindeberg} then controls $q'$ and $r'$
against the Gaussian moments with error
$C_P[u+(1+W)^{3/2}(N^{-1/2}+\|e\|_\infty)]$.
Both moment pairs are positive semidefinite and $|r'|,|r_G|\le1$.
The inequality $|\sqrt a-\sqrt b|\le\sqrt{|a-b|}$ therefore gives
\eqref{eq:nongaussian-one-step-state}.
Finally each $\xi_i$ is $K$-sub-Gaussian. A union bound for their
maximum, together with concentration of $N^{-1}\sum_i\xi_i^2$
around one, proves \eqref{eq:nongaussian-delocalization}.
\end{proof}

\begin{lemma}[Global stability and polynomial finite width]
\label{lem:nongaussian-global-transfer}
The map $\mathcal T_W$ in
\eqref{eq:nongaussian-gaussian-state} is globally $\sqrt W$-Lipschitz.
Suppose $H\ge2$ and a Gaussian covariance construction uses controls
$W_2,\ldots,W_H\ge1$ satisfying
\begin{equation}
 \max_\ell W_\ell\le e^{C_dD},\qquad
 \prod_{\ell=2}^H\sqrt{W_\ell}\le C_d(H+1)e^{C_dD}.
 \label{eq:control-product-bounds}
\end{equation}
Use the same controls with the rank-one transports above, starting
from the Gaussian first-layer comparison and its bounded-density
feature encoder. For every fixed accuracy $\varepsilon>0$ and
readout gain $\mathcal B$, there are constants $A,a$ such that
\begin{equation}
 N\ge A(H+1)^a e^{aD}
 \label{eq:nongaussian-polynomial-width}
\end{equation}
ensures the same bounded deterministic response is approximated in
pointwise mean squared error at most $\varepsilon^2$, uniformly in
the input and $0<\nu\le1$. Constants may depend on $P,d,\varepsilon$
and $\mathcal B$, but not on $H,D,N,\nu$.
\end{lemma}
\begin{proof}
The map from $Y\in L^2$ to
$(\E\xi_GY,\|Y-\xi_G\E\xi_GY\|_2)$ is nonexpansive: orthogonal
projection controls the first coordinate, and the reverse triangle
inequality controls the residual norm. Since $\phi$ is
$1$-Lipschitz and $\nu\le W$, this proves the global Lipschitz claim.

For clarity, the first direction is formed from
$a_i=\psi(W_i^1)/(\mathcal B L_H\sqrt N)$ and an independent uniform
unit vector in $a^\perp$, where the scalar encoder is bounded by
$Q$ and $L_H\ge4Q/\mathcal B$. Then $\|a\|\le1/4$ deterministically;
completing to a unit vector preserves infinity norm
$O(\sqrt{\log(N/\eta)/N})$ with failure at most $\eta$.
The first-layer empirical $q,r$ moments have the same concentration
as in the Gaussian construction. This is a comparison-law operation,
not conditioning the original prior on a support event.

Apply Lemma~\ref{lem:nongaussian-state-layer} at each layer, allocating
failure probability $\eta/H$. Every new direction satisfies the
same infinity-norm bound. Take
$u=C_P\sqrt{\log(2H/\eta)/N}$.
The one-step state error is at most
\[
 C_P(1+\max W_\ell)
       [\log(2NH/\eta)/N]^{1/4}.
\]
Global Lipschitz propagation bounds the total error by this quantity
times $H\prod\sqrt{W_\ell}$, plus the first-layer error with the
same product. The polynomial bounds
\eqref{eq:control-product-bounds} imply that
\eqref{eq:nongaussian-polynomial-width}, with a sufficiently large
fixed exponent, makes the final state error as small as required.
The logarithm of $N$ is absorbed by increasing that exponent.
This analysis is pointwise in the input, with constants uniform over
all inputs; it therefore bounds the supremum of pointwise expected
squared errors without taking an input-space union bound.

The final readout mean is $\mathcal B e_H/\sqrt N$ and its native
centered noise contributes at most $1/N$ to pointwise output MSE.
Both the finite-width mean response and the limiting response have
absolute value at most $\mathcal B$. Choose the total exceptional
probability $\eta$ small relative to $\varepsilon^2/\mathcal B^2$,
and then the on-event state error sufficiently small. This completes
the MSE bound. Auxiliary directions and adaptive transports are
measurable functions of previous weights and independent auxiliary
randomness; entropy is bounded by the conditional chain rule and
convexity under removal of the auxiliaries.
\end{proof}
